\chapter{The K4510 and the Commander X16}\label{cha:x16}

The Commander X16 is the other new 8-bit machine of this
decade: a real board, a 65C02 at 8\,MHz, VERA for video and
PSG sound, a YM2151, and a BASIC that is Commodore's V2 with
the X16's words added. The K4510 is a machine that never
existed, with a 45GS02 at 40.5\,MHz, VICKY, MELODY and a
BASIC that compiles. They were never meant to be compared,
which is exactly why it is instructive to do so.

This appendix runs the same small programs on both. The
programs are in the repo under \pth{compare/x16/}; the K4510
halves are also on the machine in \pth{/LANG/BASIC/EX} and
\pth{/LANG/BASIC/EX/X16}, where you can run them yourself.
One script, \pth{compare/x16/run.py}, runs every program on
both emulators headless and writes the table.

\begin{aside}
\textbf{Read the caveats first.} Every number here is honest
and most of them are unfair, in the way all cross-machine
numbers are. K4510 BASIC is compiled; X16 BASIC is
interpreted. The K4510's clock is five times the X16's. VICKY
has a blitter; VERA does not. A ratio of fifty to one in the
first table says nothing about the two CPUs: it says what a
person typing the same program into each machine gets back.
The one place the CPUs meet on equal terms is the second
table's 8\,MHz column.
\end{aside}

\section{How the clocks are read}
Time is always the machine's own clock, never the host's. On
the X16 it is \texttt{TI}, the jiffy counter at 60\,Hz
(\texttt{cc65}'s \texttt{clock()} reads the same). On the
K4510 it is \texttt{FRAMES} in BASIC and the frame counter at
\texttt{SYS+\$0D} in C, also 60\,Hz. Neither emulator is
hurried: \texttt{x16emu} never runs with \texttt{-warp} while
timing, and the K4510's headless harness runs as fast as the
host allows but counts frames, not seconds. The resolution is
a sixtieth of a second on both, which is why the K4510's
Rugg/Feldman programs run their loop a hundred times and
divide: at 40\,MHz, compiled, one run is over before the
counter moves.

\section{Track 1: each machine's own BASIC}
The Rugg/Feldman benchmarks (Kilobaud and \emph{Personal
Computer World}, 1977), the Byte Sieve (1981), Gordon
Henderson's text Mandelbrot and a string-handling program,
each written as a user of that machine would write it. The
X16 listings are numbered and use \texttt{GOTO}; the K4510
ones use \texttt{DO}/\texttt{LOOP} and labels, except
Mandelbrot, which is Henderson's listing untouched on both
(his benchmark asks for that).

Benchmark 1, on the X16 and then the K4510:

\begin{type}
100 REM RUGG/FELDMAN BENCHMARK 1 (1977) -- X16 BASIC
130 PRINT "RUGG/FELDMAN BENCHMARK 1, X16 BASIC"
200 T=TI
300 FOR K=1 TO 1000
310 NEXT K
900 T=TI-T
910 PRINT "RESULT RF1";T/60;"S"
920 PRINT "DONE"
\end{type}

\begin{type}
' Rugg/Feldman benchmark 1 (Kilobaud/PCW, 1977) -- K4510 BASIC
CONST RUNS = 100
t = PEEK(54541) + 256 * PEEK(54542)   ' the 60 Hz frame counter
FOR r% = 1 TO RUNS
  FOR k = 1 TO 1000
  NEXT k
NEXT r%
t = PEEK(54541) + 256 * PEEK(54542) - t
s = t / 60 / RUNS
PRINT "THIS MACHINE, ONE RUN: "; STR$(INT(s * 10000 + 0.5) / 10000); " S"
\end{type}

The Sieve keeps its 8191 flags in an integer array on both
machines: a float array that size (40\,KB) does not fit in the
X16's 38\,KB of BASIC memory, and CBM BASIC converts every
integer-array element to a float and back on each access.
That is the X16's cost to bear, and it is what a user would
write. The string program builds a 200-character string one
\texttt{CHR\$} at a time, reverses it with \texttt{MID\$},
counts its vowels, does 200 \texttt{STR\$}/\texttt{VAL} round
trips and 100 \texttt{LEFT\$}/\texttt{RIGHT\$} comparisons,
five times over; on the X16 the interpreter's garbage
collector runs inside that time, as it would for anyone.

\begin{center}
\small
\begin{tabular}{@{}l l r r r@{}}
\toprule
\textbf{Program} & \textbf{What} & \textbf{X16 (s)} & \textbf{K4510 (s)} & \textbf{X16 / K4510} \\
\midrule
RF1 & 1000 empty FOR/NEXT & 0.18 & 0.0080 & 23x \\
RF2 & 1000 turns, K=K+1 / IF\ldots GOTO & 1.20 & 0.0077 & 156x \\
RF3 & + A=K/K*K+K-K & 2.22 & 0.0307 & 72x \\
RF4 & + A=K/2*3+4-5 & 2.43 & 0.0298 & 82x \\
RF5 & + GOSUB & 2.67 & 0.0317 & 84x \\
RF6 & + FOR L=1 TO 5 & 4.10 & 0.0758 & 54x \\
RF7 & + M(L)=A & 6.55 & 0.159 & 41x \\
RF8 & 100 $\times$ K\^{}2, LOG, SIN & 1.00 & 0.0067 & 149x \\
SIEVE & Byte Sieve, 8191 flags, once & 39.40 & 0.32 & 124x \\
MANDEL & Henderson's text Mandelbrot & 118.08 & 2.25 & 52x \\
STRINGS & build, reverse, scan, STR\$/VAL & 8.38 & 7.88 & 1.1x \\
\bottomrule
\end{tabular}
\end{center}
\noindent\emph{The 2026-10-09 run: x16emu r49, cc65 2.19, both at their
own 60\,Hz clocks. For scale, the 1977 figures for a C64 are 1.2, 9.3,
17.6, 19.5, 21.0, 29.5, 47.5 and 11.5 seconds: the X16 is the C64's
BASIC at eight times the clock, and the table says so.}

Is there an interpreted BASIC on the K4510 that would make a
fairer pair? Not now. EhBASIC, the machine's first BASIC, was
retired in October 2026 when K4510 BASIC replaced it. BBC
BASIC runs on the Tube co-processor, which is a host process
at host speed, so its times mean nothing against a machine
clock. RX is an interpreter, but it is not BASIC. So the
first table stays what it is: the BASIC each machine hands
you.

\section{Track 2: the same C on both}
Three C programs, one source each, built by the same
\texttt{cc65} for both machines: \texttt{cl65 -t cx16 -O} for
the X16, and for the K4510 the recipe \texttt{CC} uses at the
prompt (\texttt{cc65 -t none --cpu 65c02 -O}, linked with
\pth{demo/prg.cfg}). The sieve again (ten iterations, as
Gilbreath ran it), a tight 16-bit integer loop of 200\,000
turns, and a memory copy: 256 \texttt{memcpy()}s of 4\,KB by
the C library (the same generic routine in both libraries),
then 64 copies of the same block by a plain byte loop. A small
header, \pth{compare/x16/c/bench.h}, hides the only
differences: where the clock is, and how a character is
printed.

The K4510 runs each program twice: at its own 40.5\,MHz, and
again with the harness told \texttt{K4510\_CPU\_HZ=8000000},
the X16's clock. The settings menu's own ladder stops at
10\,MHz; 8 is a harness setting for this purpose. In that
column the two CPUs run the same 65C02 instructions at the
same speed, and nothing in the code touches the 45GS02's
wider registers.

\begin{type}
    for (j = 0; j < 200; j++)
        for (i = 0; i < 1000; i++)
            acc += (i ^ (acc >> 3)) & 0xFF;
\end{type}

\begin{center}
\small
\begin{tabular}{@{}l r r r r r@{}}
\toprule
\textbf{Program} & \textbf{X16 8\,MHz} & \textbf{K4510 40.5} & \textbf{K4510 8\,MHz} & \textbf{X16 / 40.5} & \textbf{X16 / 8} \\
\midrule
CSIEVE (10 $\times$ 8191) & 5.57 & 0.88 & 4.43 & 6.3x & 1.3x \\
CLOOP (200\,000 turns) & 5.83 & 0.87 & 4.35 & 6.7x & 1.3x \\
CMEMCPY (1 MB) & 2.00 & 0.32 & 1.65 & 6.2x & 1.2x \\
CBYTECOPY (256 KB) & 4.53 & 0.72 & 3.65 & 6.3x & 1.2x \\
\bottomrule
\end{tabular}
\end{center}
\noindent\emph{Seconds. The same cc65 output on both; the K4510 twice.}

\section{Track 3: graphics and sound}
Each machine by its natural means, from BASIC. The X16's
\texttt{SCREEN 128} is 320$\times$240 in 256 colours;
\texttt{RECT}, \texttt{LINE} and \texttt{PSET} draw on it
through the KERNAL's graphics routines, a pixel at a time
through VERA's data port. The K4510's \texttt{GRAPHICS 1} is
the same 320$\times$240 picture laid over the text;
\texttt{BOX} and \texttt{LINE} are one blitter operation each.
Sprites: one 16$\times$16 shape at 8 bits a pixel, poked into
VRAM (\texttt{VPOKE}) or far memory (\texttt{POKE} above
65535), sixteen sprites sharing it, a hundred moves of all
sixteen. On both machines that measures how fast BASIC can
issue 1600 \texttt{MOVSPR}s, which is not the same as how fast
the hardware can move pixels; the K4510's \texttt{SPRVEL}
would let VICKY do it alone.

\begin{type}
' K4510: a shape in far memory, sixteen sprites on it
CONST SHAPE = 262144                  ' $040000, sprite page 1024
FOR i% = 0 TO 255: POKE SHAPE + i%, 16 + (i% AND 15): NEXT
FOR n% = 0 TO 15: SPRDEF n%, 1024, 16, 16, 8: NEXT
\end{type}

\begin{type}
140 REM X16: THE SAME SHAPE IN VRAM AT $1:3000
140 FOR I=0 TO 255:VPOKE 1,12288+I,16+(I AND 15):NEXT
160 SPRMEM N,1,12288,1
170 SPRITE N,3,0,0,1,1,1
\end{type}

The chord is not timed, because a chord is not a race. The
X16 plays C~major on three YM2151 channels with
\texttt{FMCHORD}, then on the VERA PSG with
\texttt{PSGCHORD}; the K4510 queues a whole note on each of
three sequencer channels with \texttt{PLAY 1,~"O4 C1"},
\texttt{PLAY 2,~"O4 E1"}, \texttt{PLAY 3,~"O4 G1"}, and MELODY
sounds them together while the program goes on. The run
script records the X16's chord to a \texttt{.wav} and checks
that it is not silence; the K4510's headless harness has no
audio device, so its chord is checked by running.

\begin{center}
\small
\begin{tabular}{@{}l l r r r@{}}
\toprule
\textbf{Program} & \textbf{What} & \textbf{X16 (s)} & \textbf{K4510 (s)} & \textbf{X16 / K4510} \\
\midrule
FILL & 200 full-screen fills, 320$\times$240 & 11.85 & 0.017 & 711x \\
LINES & 500 lines & 6.77 & 0.033 & 203x \\
POINTS & 5000 points & 9.97 & 0.35 & 28x \\
SPRITES & 16 sprites $\times$ 100 MOVSPRs & 4.83 & 0.13 & 36x \\
CHORD & C major, three voices & played & played & -- \\
\bottomrule
\end{tabular}
\end{center}

\section{What the numbers say}
\begin{itemize}
\item \textbf{Track 1 is a compiler against an interpreter, and it
  shows.} Forty to a hundred and fifty times on the loop benchmarks;
  124 times on the Sieve; 52 on Mandelbrot. The X16 figures are
  themselves about eight times a C64's, which is exactly its clock:
  X16 BASIC is the C64's BASIC, well run.
\item \textbf{The X16 holds its own on strings.} STRINGS is 8.4
  seconds against 7.9: a near draw, and the one place the compiled
  BASIC's advantage all but vanishes. The reason is where K4510 BASIC
  keeps its text: every string value is a 256-byte slot in far memory,
  copied by DMA, so each \texttt{+}, \texttt{MID\$} and
  \texttt{LEFT\$} is a block copy and a runtime call rather than a few
  instructions. CBM BASIC's strings are short, in place, and its
  garbage collector, slow as it is, was not asked to do much here. If
  the K4510 has a BASIC to improve, this is where.
\item \textbf{Clock for clock, the 45GS02 is 20--30\% quicker than
  the 65C02 on the same code.} At 8\,MHz the K4510 runs the C
  programs in 1.2--1.3 times less time than the X16. That is the
  4510 family's shorter cycle counts (fewer dead cycles per
  instruction), as emulated by Xemu's core on one side and x16emu's
  65C02 on the other. The remaining 5x of the 40.5\,MHz column is the
  clock, nothing else.
\item \textbf{A blitter is a different kind of thing.} FILL at 711x
  and LINES at 203x are VICKY's blitter against the X16 KERNAL
  drawing through VERA's data port one pixel at a time, with an
  interpreter issuing the calls. POINTS at 28x and SPRITES at 36x are
  closer to the BASIC ratio, because a point or a sprite move is one
  register write on either machine: there the BASICs are being
  compared again, not the chips.
\item \textbf{Twenty fills took no frame at all} on the K4510, which
  is why the program does two hundred. Below a sixtieth of a second
  the frame counter cannot see; the Rugg/Feldman programs run their
  loop a hundred times for the same reason.
\item \textbf{Both machines got the same answers}: 1899 primes, the
  same STRINGS checksum (200, 50, 60300, 88), the same loop
  check (48952). A benchmark that disagrees with itself measures
  nothing.
\end{itemize}

\section{Running it yourself}
On a machine with the K4510 checkout, \texttt{cc65} with its
\texttt{cx16} target, and an \texttt{x16emu} with its
\texttt{rom.bin} (the official r49 release zip, unpacked, is
what the tables above used):

\begin{type}
compare/x16/run.py --x16 /path/to/x16emu-folder
compare/x16/run.py --only rf1,sieve,csieve
\end{type}

It writes \pth{compare/x16/RESULTS.md} and leaves the raw
output, the binaries, the screenshots and the chord under
\pth{compare/x16/build/}. On the K4510 itself, \texttt{CD
/LANG/BASIC/EX/X16} and \texttt{RUN FILL} (and so on; the
\texttt{RUN} matters for \texttt{FILL}, which is otherwise the
shell's memory fill). The Rugg/Feldman programs, the Sieve and
Mandelbrot (\texttt{DROGON}) are one directory up, where they
have always been.
